\documentclass{article}
\usepackage{mathtools} 
\usepackage{fontspec}
\usepackage[UTF8]{ctex}
\usepackage{amsthm}
\usepackage{mdframed}
\usepackage{xcolor}
\usepackage{amssymb}
\usepackage{amsmath}
\usepackage{hyperref}


% 定义新的带灰色背景的说明环境 zremark
\newmdtheoremenv[
  backgroundcolor=gray!10,
  % 边框与背景一致，边框线会消失
  linecolor=gray!10
]{zremark}{注释}

% 通用矩阵命令: \flexmatrix{矩阵名}{元素符号}{行数}{列数}
\newcommand{\flexmatrix}[4]{
  \[
  #1 = \begin{pmatrix}
    #2_{11}     & #2_{12}     & \cdots & #2_{1#4}   \\
    #2_{21}     & #2_{22}     & \cdots & #2_{2#4}   \\
    \vdots      & \vdots      & \ddots & \vdots     \\
    #2_{#31}    & #2_{#32}    & \cdots & #2_{#3#4}
  \end{pmatrix}
  \]
}

% 简化版命令（默认矩阵名为A，元素符号为a）: \quickmatrix{行数}{列数}
\newcommand{\quickmatrix}[2]{\flexmatrix{A}{a}{#1}{#2}}


\begin{document}
\title{5.2注释}
\author{张志聪}
\maketitle

\begin{zremark}
  关于命题2.2 的注意点；
\end{zremark}

\begin{itemize}
  \item 在证明必要性时，按照题设“使$f$变为$g$”其实隐含了$f = g$。

        否则命题在证明时，无法使用命题2.1；

  \item 在思考二次型问题时，不能按照数学分析的思考方式，来确定$f$和$g$的相等问题（即：函数相等的问题）。

        对于线性映射来说，换了一组基，二次型就会改变，
        如果以$X = Y \in K^n$思考，$f = g$无法成立，
        而是应该考虑成
        \begin{align*}
          \alpha = X = T Y
        \end{align*}
\end{itemize}

\end{document}
